% Paper 6: Carbon-ActiveAgent — Active Inference for HLS DSE
% Target: DAC / NeurIPS AI4Science Workshop

\documentclass[conference]{IEEEtran}
\usepackage{amsmath,graphicx,booktabs,hyperref}

\title{Active Inference for HLS Design Space Exploration:\\
When Should an LLM Agent Learn from Synthesis Feedback?}

\author{
\IEEEauthorblockN{William Chen}
\IEEEauthorblockA{National Taiwan University}
}

\begin{document}
\maketitle

\begin{abstract}
We investigate whether LLM agents can improve HLS pragma predictions
through iterative synthesis feedback. Combining ideas from Active
Inference (prediction-error-driven learning) and tool-augmented agents,
we build Carbon-ActiveAgent: an agent that predicts synthesis outcomes,
observes actual Vitis HLS results, and updates a hierarchical memory
when surprised. Our experiments reveal three counter-intuitive findings:
(1) the first LLM prediction (R0) is consistently the best---iterative
refinement with feedback does not improve over single-shot prediction;
(2) injecting learned rules into the prompt \emph{increases} failure
rate from 7\% to 40\% (Context Distraction); and (3) fine-tuning a 2B
model on 3,409 ForgeHLS samples produces PFS=0 despite 60\% loss
reduction, suggesting a minimum model capacity threshold. These
negative results establish important boundaries for LLM-based hardware
optimization and motivate future work on larger fine-tuned models
with synthesis-reward training.
\end{abstract}

\section{Introduction}

Traditional HLS Design Space Exploration (DSE) tools like
AutoDSE~\cite{autodse} search the pragma space through hundreds
of synthesis iterations. Recent LLM-based approaches offer single-shot
pragma prediction but lack the ability to improve from feedback.

We ask: \emph{Can an LLM agent learn from synthesis results to
progressively improve its predictions?}

Inspired by Active Inference~\cite{friston2010}---the theory that
biological agents minimize prediction error through a
predict-observe-update loop---we design Carbon-ActiveAgent with three
components:

\begin{enumerate}
    \item \textbf{Memory Pyramid}: Hierarchical storage (L0: events,
    L1: episode digests, L2: learned rules) enabling cross-kernel
    knowledge transfer.
    \item \textbf{Surprise Gate}: Measures prediction error between
    expected and actual synthesis metrics, triggering learning only
    when surprised.
    \item \textbf{Iterative Refinement}: LLM generates pragmas $\to$
    Vitis synthesizes $\to$ structured feedback $\to$ LLM refines.
\end{enumerate}

\section{Architecture}

\subsection{Memory Pyramid}

The memory pyramid compresses optimization experience into three
levels, inspired by ContinuousAgent~\cite{continuous}:

\begin{itemize}
    \item \textbf{L0 (Working)}: Raw synthesis results from current
    optimization round. Discarded after L1 compression.
    \item \textbf{L1 (Episodes)}: Digests of past kernel optimizations
    (e.g., ``FIR: PFS=0.667, 3 pragmas''). Keeps last 20 episodes.
    \item \textbf{L2 (Rules)}: Cross-kernel rules learned from
    surprising outcomes (e.g., ``MAC-heavy kernels benefit from
    coefficient array partitioning'').
\end{itemize}

This enables constant-size context regardless of optimization history.

\subsection{Surprise Gate}

Before each synthesis, the agent predicts expected improvement based
on pragma types (PIPELINE, UNROLL, PARTITION). After synthesis, the
prediction error determines learning intensity:

\begin{equation}
    \text{surprise} = |\text{improvement}_{\text{predicted}} -
    \text{improvement}_{\text{actual}}|
\end{equation}

If surprise exceeds threshold $\tau$, the outcome is analyzed and
a new L2 rule is generated. Low-surprise outcomes are expected and
require no memory update.

\subsection{Structured Feedback}

Vitis HLS reports are converted to structured feedback with explicit
resource budgets:

\begin{verbatim}
Latency: 1028 cycles
DSP: 6/50 (12%) [OK]
LUT: 1672/10000 (17%) [OK]
Rules: Change at most 2 pragmas per round
\end{verbatim}

This prevents the ``DSP explosion'' observed in unstructured feedback
(where the model interpreted ``you have room'' as ``use maximum
resources'').

\section{Experiments}

\subsection{Closed-Loop Iterative Refinement}

We run 3 kernels $\times$ 5 rounds on Gemma4, with Vitis HLS
synthesis on AWS between each round (Table~\ref{tab:iterative}).

\begin{table}[t]
\centering
\caption{Closed-loop iterative refinement. R0 (first prediction)
is consistently the best. Subsequent rounds oscillate or regress.}
\label{tab:iterative}
\begin{tabular}{lccccc}
\toprule
\textbf{Kernel} & \textbf{R0} & \textbf{R1} & \textbf{R2} & \textbf{R3} & \textbf{Best} \\
\midrule
FIR      & \textbf{18}  & 64    & 2050  & 18   & \textbf{18} \\
DCT      & \textbf{34}  & 514   & 258   & 258  & \textbf{34} \\
DotProd  & \textbf{257} & 257   & FAIL  & 258  & \textbf{257} \\
\bottomrule
\end{tabular}
\end{table}

\textbf{Finding 1}: R0 is the ceiling. The model's best pragma
suggestion occurs at the first attempt. Feedback-based refinement
cannot surpass single-shot prediction for 4.5B models.

\subsection{Memory Injection A/B Test}

We compare baseline (no memory) vs. memory-augmented prompts
on 5 kernels $\times$ 3 seeds (Table~\ref{tab:memory}).

\begin{table}[t]
\centering
\caption{Memory injection A/B test. Memory context does not
significantly improve PFS and increases failure rate.}
\label{tab:memory}
\begin{tabular}{lcccc}
\toprule
\textbf{Condition} & \textbf{Mean PFS} & \textbf{Std} & \textbf{Fail\%} & \textbf{t-stat} \\
\midrule
Baseline (no memory)  & 0.889 & 0.097 & 7\%  & --- \\
With memory context   & 0.933 & 0.070 & 40\% & 1.25 \\
\bottomrule
\end{tabular}
\end{table}

\textbf{Finding 2}: Memory injection is not statistically significant
($t=1.25 < 2.0$) and \emph{increases} failure rate from 7\% to 40\%.
This is another manifestation of Context Distraction: learned rules
injected as prompt context compete with the target code for the
model's limited attention.

\subsection{Fine-tuning Attempt}

We LoRA fine-tune gemma-2-2b on 3,409 ForgeHLS Pareto samples
(Table~\ref{tab:finetune}).

\begin{table}[t]
\centering
\caption{Fine-tuning results. Loss decreases but PFS remains zero,
suggesting a model capacity threshold.}
\label{tab:finetune}
\begin{tabular}{lcccc}
\toprule
\textbf{Config} & \textbf{Steps} & \textbf{Loss} & \textbf{PFS} & \textbf{Output} \\
\midrule
50 steps, full code    & 50   & 19$\to$13.8 & 0.000 & Garbled \\
500 steps, full code   & 500  & 19$\to$7.3  & 0.000 & PIPELINE OFF \\
3 epochs, pragma-only  & 626  & 19$\to$7.6  & 0.000 & Repeats prompt \\
\midrule
\emph{gemma4 zero-shot} & 0   & ---         & \textbf{0.645} & Valid pragmas \\
\bottomrule
\end{tabular}
\end{table}

\textbf{Finding 3}: Despite 60\% loss reduction, the 2B model
produces no valid pragma output. Meanwhile, the 4.5B MoE model
achieves PFS=0.645 \emph{without any fine-tuning}. This suggests
a minimum model capacity threshold for HLS pragma prediction that
gemma-2-2b (2B dense) falls below, regardless of training data.

\section{Discussion}

\subsection{Why R0 is the Ceiling}

Small models (4.5B) appear to express their full hardware optimization
knowledge in a single inference. Additional context---whether from
synthesis feedback, memory rules, or retrieved examples---acts as
noise rather than signal. This is consistent with our Context
Distraction finding (Paper 2).

\subsection{The Knowledge Encoding Problem}

Three approaches to encoding HLS knowledge were tested:
\begin{enumerate}
    \item \textbf{Prompt injection} (memory, RAG): Hurts due to
    Context Distraction.
    \item \textbf{Fine-tuning} (LoRA): Fails on 2B model; needs
    larger base.
    \item \textbf{Zero-shot prompting}: Works best for 4.5B MoE.
\end{enumerate}

The most effective knowledge encoding for small models is
\emph{architectural}---MoE expert routing naturally specializes
for different pragma types without explicit training.

\subsection{Implications for AutoDSE Comparison}

AutoDSE achieves near-Pareto results through 100--500 synthesis
iterations. Our agent achieves PFS=0.645 in one shot but cannot
improve further. The two approaches are complementary:
CarbonCode provides a strong warm-start, reducing AutoDSE's
search space by eliminating clearly suboptimal configurations.

\subsection{Future Directions}

\begin{enumerate}
    \item \textbf{Larger model fine-tuning}: LoRA on 7B+ models
    with ForgeHLS data may cross the capacity threshold.
    \item \textbf{DPO training}: Direct Preference Optimization
    using (Pareto, worst-case) code pairs as preference signal.
    \item \textbf{Synthesis reward}: Using Vitis latency as reward
    signal for RLHF, enabling true closed-loop learning.
\end{enumerate}

\section{Conclusion}

We present Carbon-ActiveAgent, an Active Inference framework for
HLS optimization. Our honest negative results---iterative refinement
does not improve R0, memory injection increases failure rate,
fine-tuning 2B models yields PFS=0---establish important boundaries
for LLM-based hardware design. The most effective approach for
current small models remains simple zero-shot prompting with
post-process merge. Crossing these boundaries requires either
larger models or fundamentally different training approaches
(DPO, synthesis-reward RL).

\bibliographystyle{IEEEtran}
\begin{thebibliography}{3}
\bibitem{autodse} A. Sohrabizadeh et al., ``AutoDSE,'' ACM TODAES, 2022.
\bibitem{friston2010} K. Friston, ``The free-energy principle,'' Nature Rev. Neuroscience, 2010.
\bibitem{continuous} ContinuousAgent, ``Active Inference for Autonomous Agents,'' 2025.
\end{thebibliography}

\end{document}
